\documentclass[10pt,a4paper]{amsart}
\usepackage[margin=19mm]{geometry}
\usepackage{amsmath,amssymb,mathtools}
\usepackage[scr=rsfs,cal=euler]{mathalfa}
\usepackage{xfrac}
\usepackage[unicode,hidelinks,bookmarks=false]{hyperref}
\newcommand{\ie}{i.e.\ }
\newcommand{\cf}{cf.\ }
\newcommand{\resp}{respectively\ }
\newcommand{\Subsection}{\subsection}
\providecommand{\nobreakdash}{\nobreak\hbox{-}}
\setlength{\emergencystretch}{2em}
\multlinegap0pt
\usepackage{bm}
\usepackage{bbm}
\usepackage{dsfont}
\usepackage{xspace}
\usepackage{cancel}
\usepackage{mathtools}
\usepackage{amsthm}
\makeatletter
\@ifundefined{thm}{\newtheorem{thm}{Thm}}{}
\@ifundefined{defn}{\newtheorem{defn}[thm]{Definition}}{}
\@ifundefined{lem}{\newtheorem{lem}[thm]{Lemma}}{}
\makeatother
\makeatletter
\DeclareMathOperator{\Aut}{Aut}
\expandafter\ifx\csname com\endcsname\relax
\newenvironment{com}%
% begin comment
{\ifthenelse{\equal{\showcomments}{yes}}%
% then begin comment in margin
{\footnotemark
        \begin{lrbox}{\commentbox}
        \begin{minipage}[t]{1.25in}\raggedright\sffamily\tiny
        \footnotemark[\arabic{footnote}]}
% else eat contents of the environment
{\begin{lrbox}{\commentbox}}}%
% end comment
{\ifthenelse{\equal{\showcomments}{yes}}%
% then end comment
{\end{minipage}\end{lrbox}\marginpar{\usebox{\commentbox}}}
% else finish eating
{\end{lrbox}}}
\fi
\newcommand{\dist}{\textup{\textsf{d}}}
\expandafter\ifx\csname enumeratecontinue\endcsname\relax
\newenvironment{enumeratecontinue}{
  \setcounter{enumitemp}{\value{enumi}}
  \begin{enumerate}
  \setcounter{enumi}{\value{enumitemp}}
}
{
  \end{enumerate}
}
\fi
\makeatother
\title[Cubulating mapping tori of some polynomial growth free group]{Cubulating mapping tori of some polynomial growth free group automorphisms}
\author{Mark Hagen, Daniel T Wise}
\date{Source: arXiv:1605.07879v3 (CC BY 4.0)}
\begin{document}
\maketitle

\noindent\emph{Source and licence.} Cubulating mapping tori of some polynomial growth free group automorphisms. Version arXiv:1605.07879v3 (CC BY 4.0), distributed under the Creative Commons Attribution 4.0 International licence, as recorded in the arXiv licence metadata for this exact version and shown on the abstract page https://arxiv.org/abs/1605.07879v3. Full rights notice: licenses/arxiv-1605.07879-CC-BY-4.0.txt.\par\smallskip

\noindent\textit{Editorial scope.} Counted unit: the Lem whose source label is \texttt{lem:star} in the article (source0.tex lines 563--572, source label \texttt{lem:star}), delivered together with the article's own complete original proof of it (source0.tex lines 574--593, one \texttt{proof} environment of 20 lines). No mathematics is written, repaired or re-derived: statement and proof are the article's own text line for line. Required context retained uncounted: defn:finite-set-elevation (lines 559--561). Every definition, display and label of those lines is retained unchanged. Omitted from this package and not redistributed: 1--558, 562--562, 573--573, 594--2062. Nothing inside a retained excerpt is omitted or altered: every retained body is delivered exactly as the article's lines are written, so no omission receipt of any kind accompanies this package. Citations: the retained excerpts carry no citation command at all, so no bibliography entry of the article is transcribed and no reference is redistributed. Reference adaptation: none was needed and none was made; every reference inside the delivered unit resolves inside it, so no unresolved reference or citation is produced. Printed slips kept literal: no printed word, symbol, hypothesis or number of the counted statement or of its proof was altered. Layout and class: the article's own class is \texttt{amsart} with packages including amsmath, ifthen, amsfonts, amssymb, srcltx, amsopn, color, enumerate, xy, pb-diagram, pb-xy, overpic, mathabx, fontenc; the wrapper is the lane's pinned \texttt{amsart} editorial wrapper at 10pt on A4, which restates the theorem-like environments the retained text needs and adds the article's own macro definitions that the retained text needs (\texttt{\textbackslash Aut}, \texttt{\textbackslash dist}), with extra wrapper packages (bm, bbm, dsfont, xspace, cancel, mathtools). The delivered numbering is the unit's own. Assets: no retained span contains a figure environment, a graphics-inclusion command, a PDF-page inclusion, a table or a TikZ picture; no image file is copied into this package, the package declares no asset dependency and no image file is redistributed with it. Licence: version arXiv:1605.07879v3 of this article is distributed under the Creative Commons Attribution 4.0 International licence, as recorded in the arXiv licence metadata for this exact version and shown on the abstract page https://arxiv.org/abs/1605.07879v3. A full rights notice is delivered at licenses/arxiv-1605.07879-CC-BY-4.0.txt. Characters: every retained excerpt is pure ASCII, so the delivered engine, XeLaTeX with its default Latin Modern text font, reports no missing character.
\begin{defn}[Algebraic elevation]\label{defn:finite-set-elevation}
Let $G$ be a group and let $G'\leq G$ have finite index.  Let $\mathcal Q\subset G$.  Let $\mathcal Q_0\subset \mathcal Q$ contain exactly one representative of each $G$--conjugacy class in $\mathcal Q$.  For each $g\in G$ and each $q\in\mathcal Q_0$, let $q'$ be a generator of $g\langle q\rangle g^{-1}\cap G'$.  Let $\mathcal Q''$ be the set of all such $q'$ (as $g$ and $q$ vary) and let $\mathcal Q'\subset G'$ contain exactly one representative of each $G'$--conjugacy class in $\mathcal Q''$.  Then $\mathcal Q'$ is an \emph{elevation} of $\mathcal Q$.
\end{defn}

\begin{lem}[Induced cubulation]\label{lem:star}
Let $G$ be a group and $G'\leq G$ a finite index subgroup.  If $G'$ acts on a CAT(0) cube complex $C'$, then $G$ acts on a CAT(0) cube complex $C$ such that:
\begin{itemize}
 \item $C$ is isomorphic to $(C')^{[G:G']}$, with the product cubical structure.
 
 \item For $a,b\in G$ with $\|h_1ah_1^{-1}\|_{C'}^{virt}=\|h_2bh_2^{-1}\|_{C'}^{virt}$ for all $h_1,h_2\in G$, we have $\|a\|_C=\|b\|_C$.
\end{itemize}

Suppose, moreover, that $\mathcal Q\subset G$ and $\mathcal Q'\subset G'$ is an elevation (see Definition \ref{defn:finite-set-elevation}).  Suppose that $\mathcal Q'$ is wall-independent for the $G'$--action on $C'$, and that for all $q,p\in\mathcal Q$ and $g\in G$, $g\langle p\rangle g^{-1}\cap \langle q\rangle\neq\{1\}$ implies $p=q^{\pm1}$. Then $\mathcal Q$ is wall-independent for the $G$--action on $C$.
\end{lem}

\begin{proof}
Let $\rho:G'\to\Aut(C')$ be the given action.  Equipping $G$ with the left multiplication action and $C'$ with the action $\rho$, let $C$ be the set of maps $f:G\to C'$ such that $f(hg)=\rho(h)(f(g))$ for all $h\in G'$ and $g\in G$.  Define a $G$--action on $C$ by: given $f\in C$ and $g\in G$, let $g\cdot f:G\to C'$ be given by $g\cdot f(k)=f(kg)$ for all $k\in G$.

Let $g_1,\ldots,g_{[G:G']}$ be a right transversal for $G'$ in $G$.  Define a map $I:C\to (C')^{[G:G']}$ by $I(f)=(f(g_1),\ldots,f(g_{[G:G']}))$.  This is bijective since $G'$--equivariant maps from $G$ are determined uniquely by where they send the $g_i$, and hence, by pulling back the cubical structure/metric, we can regard $C$ as a CAT(0) cube complex with the product structure as in the statement.


The inverse $I:(C')^{[G:G']}\to C$ is given by $I((c_i)_i)=f$, where $f(hg_i)=\rho(h)c_i$ for $h\in G'$ and $i\leq [G:G']$.  The $G$--action on $(C')^{[G:G']}$ is thus given by $g\cdot(c_i)_i=I(g\cdot f)=(f(g_ig))_i$.  So if $g=g_j^{-1}ag_j$ for some $a\in G'$, then for any $(c_i)_i$, the $j$--coordinate of $g\cdot(c_i)_i$ is $$f(ag_j)=\rho(a)f(g_j)=\rho(a)c_j.$$

Let $r=[G:G']!$ and let $a\in G$.  Then $a^r\in \bigcap_ig_i^{-1}G'g_i$, so for each $i\leq [G:G']$ there exists $a_i\in G'$ such that $a^{nr}=g_i^{-1}a_i^ng_i$ for all $n>0$.  Hence, if $f\in C$, then the $i^{th}$--coordinate of $I(a^{nr}\cdot f)$ is $a^{nr}\cdot f(g_i)=f(a_i^ng_i)=\rho(a_i)^nf(g_i)$.  Hence
$$I(a^{nr}\cdot f)=(\rho(a_1)^n(f(g_1)),\ldots,\rho(a_{[G:G']})^n(f(g_{[G:G']}))).$$
Hence the $\ell_1$--distance from $I(a^{nr}\cdot f)$ to $I(f)$ is 
$\sum_i\dist_{C'}(\rho(g_ia^{r}g_i^{-1})^n(f(g_i)),f(g_i)),$
where we have used $g_ia^rg_i^{-1}=a_i$. Dividing by $n$ and passing to the limit as $n\to\infty$ gives $\|a^r\|_C=\sum_i\|g_ia^rg_i^{-1}\|_{C'}=r\sum_i\|g_iag_i^{-1}\|^{virt}_C.$  This implies that $\|a\|_C=\|b\|_C$ provided all conjugates of $a,b$ have equal virtual translation length.

\textbf{Description of walls:}  Each hyperplane of $C$ is associated to a pair $(W,g_i)$, where $W\subset C'$ is a hyperplane: $(W,g_i)$ is the set of points in $\prod_{j=1}^{[G:G']}C'$ with $i$--coordinate in $W$.  

\textbf{Wall-independence:}  Let $\mathcal Q$ and $\mathcal Q'$ be as in the statement, and suppose that $\mathcal Q'$ is wall-independent for the $G'$--action on $C'$.  Let $\mathcal Q_0$ be a set of conjugacy class representatives for $\mathcal Q$ in $G$.  Then by re-choosing the coset representatives $g_1,\ldots,g_{[G:G']}$, we can assume that the elements of $\mathcal Q'$ have the form $g_iq^{n(q,i)}g_i^{-1}$ for $q\in\mathcal Q_0$ and $1\leq i\leq [G:G']$.  

Suppose that $\mathcal Q$ is not wall-independent for the $G$--action on $C$.  Let $p,q\in\mathcal Q$ and $g\in G$ be such that  $gpg^{-1}$ and $q$ are cut by infinitely many common walls in $C$.  So by the pigeonhole principle there exists $j$ such that infinitely many walls of the form $(W,g_j)$ cut both $gpg^{-1}$ and $q$.  These same walls cut $gp^rg^{-1}$ and $q^r$, both of which belong to $g_j^{-1}G'g_j$.  Hence $g_jg p^rg^{-1}g_j^{-1}$ and $g_j^{-1}q^rg_j$ are elements of $G'$ that are both cut by infinitely many walls for the action of $G'$ on $C'$.  Now, $g_j^{-1}q^rg_j \in a \langle q_0\rangle a^{-1}$ for some $q_0\in\mathcal Q'$ and $a\in G'$ and $(g_jg)p^r(g_jg)^{-1}\in b\langle q_1\rangle b^{-1}$ for some $b\in G'$ and $q_1\in \mathcal Q'$.  So, infinitely many hyperplanes in $C'$ cut both $aq_0a^{-1}$ and $bq_1b^{-1}$, so by the wall-independence assumption, we have $aq_0a^{-1}=bq_1^{\pm1}b^{-1}$.  Hence $g\langle p^r\rangle g^{-1}$intersect $\langle q^r\rangle$ nontrivially, so by our hypothesis about $\mathcal Q$, $p=q^{\pm1}$, so $\mathcal Q$ is wall-independent.   
\end{proof}

\end{document}
